Colorectal cancer (CRC) is a significant health issue globally, with particular prevalence in Chile. This thesis is structured into two distinct phases: the first aims to discover and analyze the CRC care process within the Southern Metropolitan Public Healthcare System (SMPHS) of Chile, while the second focuses on delivering a feasible improvement strategy for CRC care in this region.

The discovery phase utilized Process Mining, an innovative approach enabling quick identification of common pathways and performance indicators within a system. This study utilized real-world data from the SMPHS to swiftly analyze and evaluate the CRC care process. Spanning from 2017 to 2023, the dataset encompassed 2048 and 2524 patient cases. Through collaboration with the Management and Networks Department of the SMPHS, key dynamics and challenges in diagnosis and treatment were uncovered and are detailed in this study.

The improvement phase presents a microsimulation model designed to simulate the implementation of a CRC screening strategy. This strategy involves conducting random fecal immunochemical tests for the population aged between 50 and 75. Positive test results lead to a colonoscopy procedure, aimed at promoting early detection of CRC. All demographic and medical data used in the model was obtained from respected sources to ensure the validity of the model. The results suggest that as population adherence rates to the screening program increase, there is a consistent reduction in late-stage cancer cases; however, overall costs also rise.

\vfill
\noindent {\bf Keywords}:  Healthcare, Process Mining, Simulation, Optimization, Colorectal Cancer, Microsimulation, Screening.
